% =====================================================================
%  courses/algebra/indice.tex — índice detallado: álgebra preuniversitaria
% ---------------------------------------------------------------------
%  5 Unidades · 46 Temas
% =====================================================================

% --- Página de índice automático de LaTeX ----------------------------------
{
\hypersetup{linkcolor=AzulPizarra}
\tableofcontents
}
\newpage

% --- Índice temático manual (resumen visual por unidades) ------------------
\newpage
\thispagestyle{fancy}

\begin{center}
  {\Large\sffamily\bfseries\color{AzulPizarra} Índice Temático — Álgebra Preuniversitaria}\\[4pt]
  {\small\sffamily\color{GrisMedio} 5 Unidades · 46 Temas · Teoría, Ejercicios y Resueltos}
\end{center}

\vspace{0.6cm}

% --- UNIDAD I --------------------------------------------------------------
\begin{tcolorbox}[
    enhanced,
    colback   = AzulPizarra,
    colframe  = DoradoAcadémico,
    arc = 3pt, boxrule = 1pt,
    left = 8pt, right = 8pt, top = 5pt, bottom = 5pt,
    fontupper = \sffamily\bfseries\color{white}
  ]
  UNIDAD I \hfill Fundamentos Algebraicos \hfill Temas 1–10
\end{tcolorbox}

\begin{multicols}{2}
  \begin{enumerate}[label=\textcolor{AzulPizarra}{\textbf{\arabic*.}},
      leftmargin=2em, itemsep=1pt]
    \item Introducción al Álgebra
    \item Conjuntos numéricos
    \item Expresiones algebraicas
    \item Operaciones algebraicas
    \item Productos notables
    \item Cocientes notables
    \item Factorización
    \item Fracciones algebraicas
    \item Radicales algebraicos
    \item Potenciación y radicación
  \end{enumerate}
\end{multicols}

\vspace{0.4cm}

% --- UNIDAD II -------------------------------------------------------------
\begin{tcolorbox}[
    enhanced,
    colback   = AzulMedio,
    colframe  = DoradoAcadémico,
    arc = 3pt, boxrule = 1pt,
    left = 8pt, right = 8pt, top = 5pt, bottom = 5pt,
    fontupper = \sffamily\bfseries\color{white}
  ]
  UNIDAD II \hfill Ecuaciones e Inecuaciones \hfill Temas 11–20
\end{tcolorbox}

\begin{multicols}{2}
  \begin{enumerate}[label=\textcolor{AzulMedio}{\textbf{\arabic*.}},
      leftmargin=2em, itemsep=1pt, start=11]
    \item Ecuaciones de primer grado
    \item Ecuaciones de segundo grado
    \item Ecuaciones bicuadradas
    \item Ecuaciones irracionales
    \item Ecuaciones exponenciales
    \item Ecuaciones logarítmicas
    \item Sistemas de ecuaciones
    \item Inecuaciones lineales
    \item Inecuaciones cuadráticas
    \item Sistemas de inecuaciones
  \end{enumerate}
\end{multicols}

\vspace{0.4cm}

% --- UNIDAD III ------------------------------------------------------------
\begin{tcolorbox}[
    enhanced,
    colback   = AzulPizarra,
    colframe  = DoradoAcadémico,
    arc = 3pt, boxrule = 1pt,
    left = 8pt, right = 8pt, top = 5pt, bottom = 5pt,
    fontupper = \sffamily\bfseries\color{white}
  ]
  UNIDAD III \hfill Relaciones y Funciones \hfill Temas 21–32
\end{tcolorbox}

\begin{multicols}{2}
  \begin{enumerate}[label=\textcolor{AzulPizarra}{\textbf{\arabic*.}},
      leftmargin=2em, itemsep=1pt, start=21]
    \item Relaciones
    \item Funciones
    \item Función lineal
    \item Función cuadrática
    \item Función polinómica
    \item Función racional
    \item Función irracional
    \item Función exponencial
    \item Función logarítmica
    \item Funciones especiales
    \item Composición de funciones
    \item Función inversa
  \end{enumerate}
\end{multicols}

\vspace{0.4cm}

% --- UNIDAD IV -------------------------------------------------------------
\begin{tcolorbox}[
    enhanced,
    colback   = AzulMedio,
    colframe  = DoradoAcadémico,
    arc = 3pt, boxrule = 1pt,
    left = 8pt, right = 8pt, top = 5pt, bottom = 5pt,
    fontupper = \sffamily\bfseries\color{white}
  ]
  UNIDAD IV \hfill Polinomios \hfill Temas 33–40
\end{tcolorbox}

\begin{multicols}{2}
  \begin{enumerate}[label=\textcolor{AzulMedio}{\textbf{\arabic*.}},
      leftmargin=2em, itemsep=1pt, start=33]
    \item Polinomios
    \item Operaciones con polinomios
    \item División algebraica
    \item Teorema del residuo
    \item Teorema de Ruffini
    \item Teorema del factor
    \item Relaciones de Girard
    \item Ecuaciones polinómicas
  \end{enumerate}
\end{multicols}

\vspace{0.4cm}

% --- UNIDAD V --------------------------------------------------------------
\begin{tcolorbox}[
    enhanced,
    colback   = AzulPizarra,
    colframe  = DoradoAcadémico,
    arc = 3pt, boxrule = 1pt,
    left = 8pt, right = 8pt, top = 5pt, bottom = 5pt,
    fontupper = \sffamily\bfseries\color{white}
  ]
  UNIDAD V \hfill Progresiones y Binomio \hfill Temas 41–46
\end{tcolorbox}

\begin{multicols}{2}
  \begin{enumerate}[label=\textcolor{AzulPizarra}{\textbf{\arabic*.}},
      leftmargin=2em, itemsep=1pt, start=41]
    \item Sucesiones
    \item Progresión aritmética
    \item Progresión geométrica
    \item Interés compuesto
    \item Binomio de Newton
    \item Triángulo de Pascal
  \end{enumerate}
\end{multicols}

\vspace{0.6cm}

\begin{cajatip}
  \textbf{Cómo usar este libro:} cada semana consta de tres partes:
  (1)~\textbf{Teoría} con definiciones, propiedades y ejemplos resueltos;
  (2)~\textbf{Ejercicios propuestos} en dos columnas con 5 alternativas;
  (3)~\textbf{Ejercicios resueltos} con desarrollo pedagógico paso a paso.
  Se recomienda estudiar la teoría primero, intentar los ejercicios y
  verificar con los resueltos.
\end{cajatip}

\newpage